We introduced HexTiles with semantic icons, a domain-agnostic superimposition-based visualization design for multivariate geospatial data.
%
With its design, qualitative results show that HexTiles is able to more easily communicate multivariaate data with MAUP with in mind as seen by the NASA TLX survery results, albeit at small sample size.
%
Further our semi-structured interview results point to the efficacy of semantic icons along with its design also puts the visualization designer in the driver's seat for guiding user attention to specific variables.
%
Our quantitative results also showed some evidence that HexTiles allows for more accurate interaction assessment, and showed strong evidence that HexTiles is more consistent in its performance and mental load. 
%
We presented feedback from domain experts that echoed these findings.
%
We also list the limitations of HexTiles gathered through our evaluations -- though some are inherent to a superimposition-based design such as HexTiles.
%
Future work includes exploring the usability of HexTiles height for data encoding, adaptive task-based hexagonal resolution, consideration of temporal data (i.e., implementing HexTiles for spatiotemporal data and the necessary design considerations), and exploring the use of animation and further interaction to better support exploration and faithful sensemaking of multivariate geospatial data.
